\documentclass[preprint,12pt,authoryear]{elsarticle}

\usepackage[utf8]{inputenc}
\usepackage{amssymb}
\usepackage{amsmath}
\usepackage{hyperref}
\usepackage{xurl}
\usepackage{booktabs}
\usepackage{graphicx}
\usepackage{float}
\usepackage{natbib}
\usepackage{array}
\usepackage{seqsplit}
\usepackage{xr-hyper}
\externaldocument{main}
\hypersetup{hypertexnames=false}

\bibliographystyle{apalike}

\setlength{\emergencystretch}{4em}
\hfuzz=6pt
\hbadness=1100

\makeatletter
\pdfstringdefDisableCommands{%
  \def\corref#1{}%
  \def\@corref#1{}%
  \def\cnotenum{}%
}
\makeatother

\journal{Finance Research Letters}

\newcommand{\BTC}{Bitcoin (BTC-USD)}
\newcommand{\GLD}{SPDR Gold Shares (GLD)}
\newcommand{\SPX}{S\&P 500 Index}
\newcommand{\IEF}{iShares 7--10 Year Treasury Bond ETF (IEF)}

% Appendix-only: left-align the title; authors/affiliations stay centered
% (elsarticle preprint default). Drop the \hrule lines that bracket the
% (empty) abstract slot.
\makeatletter
\long\def\pprintMaketitle{\clearpage
  \iflongmktitle\if@twocolumn\let\columnwidth=\textwidth\fi\fi
  \resetTitleCounters
  \def\baselinestretch{1}%
  \printFirstPageNotes
  \thispagestyle{pprintTitle}%
  \def\baselinestretch{1}%
  \begin{flushleft}%
    \Large\@title\par
  \end{flushleft}%
  \vskip18pt
  \ifx\@elsarticlenewpageafter\newpage@after@title
    \newpage
  \fi
  \begin{center}%
    \ifdoubleblind
      \vspace*{2pc}
    \else
      \normalsize\elsauthors\par\vskip10pt
      \footnotesize\itshape\elsaddress\par\vskip18pt
    \fi
    \ifx\@elsarticlenewpageafter\newpage@after@author
      \newpage
    \fi
    \ifvoid\absbox\else\unvbox\absbox\par\vskip10pt\fi
    \ifvoid\keybox\else\unvbox\keybox\par\vskip10pt\fi
    \ifx\@elsarticlenewpageafter\newpage@after@abstract
      \newpage
    \fi
  \end{center}%
  \gdef\thefootnote{\arabic{footnote}}%
}
\makeatother

\begin{document}

\begin{frontmatter}
\sloppy

\title{{\large Appendix}\\[8pt] Regime Labels Are Not Representation-Invariant:\\ Implications for Model Risk Governance}

\author[1]{Kai Zheng}
\ead{kaizhengnz@gmail.com}
\author[2]{Rand Low}
\ead{rlow@bond.edu.au}
\author[1]{Ruili Wang\corref{cor1}}
\ead{ruili.wang@massey.ac.nz}

\cortext[cor1]{Corresponding author: Ruili Wang}

\affiliation[1]{organization={School of Mathematical and Computational Sciences\\Massey University},
            city={Auckland},
            country={New Zealand}}

\affiliation[2]{organization={Bond Business School\\Bond University},
            city={Gold Coast},
            state={QLD},
            country={Australia}}

\end{frontmatter}
\sloppy

% S-prefixed numbering for sections, tables, figures, equations
\renewcommand{\thesection}{S\arabic{section}}
\renewcommand{\thetable}{S\arabic{table}}
\renewcommand{\thefigure}{S\arabic{figure}}
\renewcommand{\theequation}{S\arabic{equation}}

The following tables and analyses constitute the Internet Appendix accompanying the main manuscript.

\section{Theoretical Background}
\label{app:theory}

\subsection{Identifiability of finite mixtures and Markov-switching models}

The theoretical foundation for state inference in regime-switching models is the identifiability of finite mixtures. \citet{Teicher1963} and \citet{YakowitzSpragins1968} establish that, under regularity conditions, the parameters of a $K$-component finite mixture are recoverable up to label permutation from sufficiently large i.i.d.\ samples. \citet{Hamilton1989} extends this to Markov-switching models with hidden state dynamics, and \citet{Guidolin2011} surveys subsequent econometric refinements. Standard asymptotic theory therefore guarantees that two analysts working with the same feature mapping $\Phi$, the same parametric family $\mathcal{M}$, and infinite data converge to the same parameters and, consequently, the same partition (up to label permutation). This guarantee underwrites essentially all backtesting-based validation of regime models in current practice.

\subsection{Limitations of the existing theory for representation invariance}
\label{app:theory_limits}

The asymptotic guarantee dissolves in three ways. First, asset returns are non-stationary on the time scales over which regime models are estimated, so the i.i.d.\ premise of \citet{Teicher1963} fails; finite-sample inference at a 252-day window is the operative regime. Second, and central to this paper, the result is conditional on a \emph{fixed} $\Phi$: there is no theorem covering two analysts who pick different $\Phi_1, \Phi_2$ on the same finite sample. The clustering-stability literature \citep{BenDavidLuxburgPal06,vonLuxburg2010} addresses stochastic perturbations of a fixed $\Phi$, not deliberate variation across economically motivated alternatives. Third, mixture identifiability presumes a known parametric emission family; Gaussian misspecification breaks the recovery guarantee even at the asymptotic limit. The supervised-learning analog is the underspecification framework of \citet{DAmouretal2022}, but regime inference lacks any held-out performance criterion to anchor the comparison. The descriptive diagnostic of Section~\ref{sec:diagnostic} is therefore the operative validation tool until a theory of representation-invariant stability exists.

\subsection{Empirical fragility of regime inference in finance}

Real-world episodes motivate this work. During the 2007--2009 GFC, regime models flagged the September 2008 transition with detection delays that varied by weeks across realised, GARCH, and implied-volatility features \citep{DanielssonJamesValenzuelaZer2016}. The 2020 COVID-19 crash compressed a regime-shift signature into one trading week, so the same model under different rolling-window horizons assigned the crisis to different states (Supplementary Figure~\ref{fig:rep_conflict}). \citet{Ampountolas2025FRL} and \citet{BlazsekKongShadoff2025} document feature- and specification-sensitivity in adjacent applications. In each case the disagreement is a systemic feature of the inference pipeline that current theoretical results do not bound, motivating an empirical diagnostic that makes the disagreement auditable.

\section{Feature Representation Definitions}
\label{app:representations}

Table~\ref{tab:rep_definitions} lists the seven admissible representations used in the baseline and step-sweep experiments. All features are computed from daily log returns $r_t = \log(P_t/P_{t-1})$ or from the adjusted-close price series $P_t$. Rolling windows are measured in business days. VaR and CVaR use a left-tail probability of $\alpha=0.05$.

% source: authors' design specification (no CSV); reps consumed by src/core/features.py
\begin{table}[!htbp]
\centering
\footnotesize
\setlength{\tabcolsep}{4pt}
\caption{Admissible feature representations.}
\label{tab:rep_definitions}
\begin{tabular}{>{\raggedright\arraybackslash}p{2.0cm}>{\raggedright\arraybackslash}p{4.5cm}>{\raggedright\arraybackslash}p{3.3cm}>{\raggedright\arraybackslash}p{1.8cm}}
\toprule
Label & Features & Windows (days) & Std. \\
\midrule
rep\_a         & Vol, Drawdown, MaxDD, VaR, CVaR   & vol=20, dd=60, tail=60    & $z$-120d \\
\seqsplit{rep\_a\_unscaled} & Vol, Drawdown, MaxDD, VaR, CVaR & vol=20, dd=60, tail=60    & None \\
rep\_b         & RealSkew, VolStab, VaR, CVaR       & skew=60, stab=60, tail=60 & $z$-120d \\
rep\_c1        & Vol, Drawdown, VaR, CVaR           & vol=30, dd=90, tail=90    & $z$-120d \\
rep\_c2        & Vol, Drawdown, MaxDD, CVaR         & vol=20, dd=60, tail=60    & $z$-120d \\
rep\_c3        & Vol, Drawdown, VaR, CVaR           & vol=20, dd=60, tail=60    & $z$-120d \\
\midrule
rep\_d         & GARCHVol, Drawdown, VaR, CVaR      & dd=60, tail=60            & $z$-120d \\
\midrule
rep\_e$^*$     & VIXLevel, VIXChange, VIXPct, CVaR  & vix\_std=120, vix\_chg=5, vix\_pct=60, tail=60 & $z$-120d \\
\bottomrule
\end{tabular}
\par\medskip
\begin{minipage}{\linewidth}
\footnotesize
\textit{Notes:} Std.\ = standardization; $z$-120d = rolling $z$-score with 120-day lookback; None = no standardization. Vol = annualised realised volatility; GARCHVol = annualised conditional volatility from a GARCH(1,1) model; Drawdown = rolling drawdown from period high; MaxDD = rolling maximum drawdown; VaR/CVaR = historical Value-at-Risk/Conditional VaR at 5\%; RealSkew = rolling realised skewness; VolStab = volatility-of-volatility. VIXLevel/VIXChange/VIXPct = rolling $z$-scored VIX level, 5-day log-change, and rolling 60-day empirical percentile of VIX. Representations rep\_a through rep\_c3 share a rolling-window architecture; rep\_d uses a parametric autoregressive volatility estimator (GARCH(1,1)); rep\_e uses the VIX implied-volatility index as a forward-looking, market-based signal. The GARCH(1,1) implementation in \texttt{src/core/features.py} falls back to a 20-day rolling realised volatility when the rolling window contains fewer than 50 observations or when the GARCH MLE fails; in the published pipeline \texttt{outputs/run.log} contains zero \texttt{GARCH\_FALLBACK} events, so rep\_d is GARCH-only and the rep\_d-vs-rep\_a contrast is not contaminated by silent fallback. $^*$rep\_e is restricted to the \SPX{} (VIX is the S\&P 500 implied vol index).

\textit{Source:} Authors' design.
\end{minipage}
\end{table}

\section{Ordering Consistency}
\label{app:ordering}

This appendix reports cross-representation ordering consistency under two metric families. The pipeline-native metric (Table~\ref{tab:ordering}) pairs states across representations via Hungarian matching on $|\Delta\mathrm{CVaR}|+|\Delta\mathrm{Vol}|$ and then asks whether the worst-CVaR state in $A$ is paired with the worst-CVaR state in $B$. We show below that this metric has a structurally inflated null that makes its near-unity observed values uninterpretable, and we report an alternative matching-free metric (Table~\ref{tab:ordering_rank_aligned}) with a null that is exactly $1/K$ by construction.

\textbf{Why the Hungarian + CVaR metric is degenerate.}\enspace
Both nulls reported in Table~\ref{tab:ordering} (permutation; independent random $K$-partitions) produce top-1 values of 0.41--0.99 across (asset, $K$), not the $1/K \in \{0.20, 0.33\}$ that a symmetric chance baseline would imply. The reason is constructive: the Hungarian optimiser minimises $|\Delta\mathrm{CVaR}|+|\Delta\mathrm{Vol}|$, so it necessarily pairs the most-extreme-CVaR state in $A$ with the most-extreme-CVaR state in $B$, and so on by rank. The top-1 statistic then asks whether this rank-aligned pairing aligns with the CVaR rank on each side, which is nearly tautological. Two independent random $K$-partitions of the same return series yield top-1 close to $1$ because their sampling-driven CVaR extremes are still matched by rank. This was stated in the previous revision of this appendix as ``expected top-1 $\approx 1/K$ under the independent null'', which empirical evidence contradicts; the claim is retracted here. Consequently the observed $0.932$ top-1 at $K=3$ is not meaningfully above its proper chance level, and the secondary finding that ``cross-representation ordering is stable even when partition agreement is not'' does not survive this correction.

\textbf{A matching-free rank-aligned metric.}\enspace
Table~\ref{tab:ordering_rank_aligned} reports a metric that does not use Hungarian matching. For each pair of representations $(A, B)$ and each rolling window, we compute per-state CVaR within each partition, relabel each partition so that label $0$ is the worst-CVaR state, label $1$ the next, and so on, and then compare the two relabelled sequences pointwise across the shared time index. Three statistics are reported: \emph{pointwise agreement} (fraction of dates on which the two relabelled sequences coincide), \emph{top-1 Jaccard} ($|A_\text{worst} \cap B_\text{worst}| / |A_\text{worst} \cup B_\text{worst}|$, where $A_\text{worst}$ is the set of dates labelled as the worst-CVaR state by $A$), and \emph{Spearman rank correlation} of the relabelled sequences. Under independent uniform random $K$-partitions the expectations are $1/K$, $1/(2K-1)$, and $0$ respectively; Monte Carlo confirms these to within $\pm 0.01$ at $n_\text{perm}=500$. The metric is computed post-hoc from the pipeline's hard-state files by \texttt{src/experiments/posthoc\_rank\_aligned\_ordering.py} (no pipeline refit).

\textbf{What Table~\ref{tab:ordering_rank_aligned} shows.}\enspace
Mean across the four assets at $K=3$: top-1 Jaccard $= 0.47$ (vs null $0.20$, a $2.4\times$ excess); pointwise agreement $= 0.55$ (vs null $0.33$, a $1.6\times$ excess); Spearman rank correlation $= 0.48$ (vs null $0.00$). The direction is consistent across $K\in\{2,3,4,5\}$: top-1 Jaccard exceeds null by factors of $1.8$--$3.4\times$ and remains genuinely above chance at higher $K$. However, the absolute values are far below the $0.93$ top-1 reported by the Hungarian metric, and are of the same order of magnitude as the partition-level cross-representation ARI in Table~\ref{tab:baseline} ($0.34$--$0.43$). Ordering structure and partition structure therefore carry comparable amounts of cross-representation information; ordering is not a rank-invariant layer that is more stable than the partition it rides on. In practical terms, when two admissible representations disagree on labels they also disagree, in roughly half of the days in question, on which state carries the worst CVaR.

% source: outputs/ordering_ci_summary.csv  (per-asset top-1 / Spearman with perm + indep nulls)
%         outputs/ordering_seed_metrics.csv (per-seed inputs)
%         outputs/ordering_null_by_k.csv    (null reference at K=3 row)
\begin{table}[!htbp]
\centering
\small
\setlength{\tabcolsep}{4.5pt}
\caption{Cross-representation ordering consistency by asset ($K=3$ baseline), with permutation and independent nulls.}
\label{tab:ordering}
\begin{tabular}{lcccccc}
\toprule
 & \multicolumn{3}{c}{Top-1 Agreement} & \multicolumn{3}{c}{Spearman Rank Corr.} \\
\cmidrule(lr){2-4}\cmidrule(lr){5-7}
Asset & Obs. & Perm Null & Indep Null & Obs. & Perm Null & Indep Null \\
\midrule
BTC-USD  & 0.946 & 0.976 & 0.932 & 0.896 & 0.930 & 0.886 \\
GLD      & 0.919 & 0.962 & 0.930 & 0.898 & 0.930 & 0.911 \\
S\&P 500 & 0.946 & 0.894 & 0.858 & 0.910 & 0.882 & 0.826 \\
IEF      & 0.919 & 0.896 & 0.848 & 0.889 & 0.892 & 0.823 \\
\midrule
Mean     & 0.932 & 0.932 & 0.892 & 0.898 & 0.909 & 0.862 \\
\bottomrule
\end{tabular}
\par\medskip
\begin{minipage}{\linewidth}
\footnotesize
\textit{Notes:} The highest-risk state in each rolling window is identified by the most negative in-state Conditional VaR (CVaR at $\alpha = 0.05$). Top-1 Agreement = fraction of windows where two representations agree on the high-risk label after Hungarian matching on $|\Delta\mathrm{CVaR}|+|\Delta\mathrm{Vol}|$. Spearman Rank Corr.\ = Spearman correlation of CVaR-based risk ranks after the same matching. Observed values are averaged across GMM and HMM models, 20 seeds, and all representation pairs. Perm Null = per-asset mean under 500 random label permutations with Hungarian matching (structurally inflated; see methodological caveat). Indep Null = both compared sequences are independent uniform random $K$-partitions, providing a symmetric chance baseline.

\textit{Source:} Authors' calculations.
\end{minipage}
\end{table}

% source: outputs/ordering_null_by_k.csv  (per-K observed vs perm/indep nulls)
\begin{table}[!htbp]
\centering
\small
\setlength{\tabcolsep}{5pt}
\caption{Cross-representation ordering consistency across $K$ (mean across assets, GMM and HMM), observed vs.\ permutation and independent null.}
\label{tab:ordering_extended}
\begin{tabular}{rcccccc}
\toprule
 & \multicolumn{3}{c}{Top-1 Agreement} & \multicolumn{3}{c}{Spearman Rank Corr.} \\
\cmidrule(lr){2-4}\cmidrule(lr){5-7}
$K$ & Observed & Perm Null & Indep Null & Observed & Perm Null & Indep Null \\
\midrule
3 & 0.932 & 0.932 & 0.892 & 0.898 & 0.909 & 0.862 \\
4 & 0.909 & 0.895 & 0.869 & 0.892 & 0.880 & 0.879 \\
5 & 0.891 & 0.678 & 0.859 & 0.896 & 0.883 & 0.890 \\
\bottomrule
\end{tabular}
\par\medskip
\begin{minipage}{\linewidth}
\footnotesize
\textit{Notes:} $K=3$ values from Table~\ref{tab:ordering}.  $K\in\{4,5\}$ observed values are means across 4 assets and 2 models (GMM, HMM).  Perm Null: per-asset mean under 500 random label permutations with Hungarian matching (structurally inflated).  Indep Null: both compared sequences are independent uniform random $K$-partitions; note that across (asset, $K$) the two null channels span $0.41$--$0.99$ for top-1 and $0.80$--$0.95$ for Spearman, far above the $1/K$ implied by a symmetric chance baseline, because the Hungarian + CVaR matching rule used to compute top-1 and Spearman constructively rank-aligns state extremes regardless of partition structure (see main text of this appendix).  Results saved to \texttt{ordering\_null\_by\_k.csv}.

\textit{Source:} Authors' calculations.
\end{minipage}
\end{table}

% source: outputs/rank_aligned_ordering_summary.csv  (per-K observed)
%         outputs/rank_aligned_ordering_pairs.csv    (per-pair detail)
\begin{table}[!htbp]
\centering
\small
\setlength{\tabcolsep}{4.5pt}
\caption{Cross-representation ordering consistency under the matching-free rank-aligned metric: relabel each partition by within-state CVaR rank (0 = worst) and compare pointwise. Null expectations are exactly $1/K$ (pointwise), $1/(2K-1)$ (top-1 Jaccard), and $0$ (Spearman) by construction.}
\label{tab:ordering_rank_aligned}
\begin{tabular}{rcccccc}
\toprule
 & \multicolumn{2}{c}{Pointwise Agreement} & \multicolumn{2}{c}{Top-1 Jaccard} & \multicolumn{2}{c}{Spearman Rank Corr.} \\
\cmidrule(lr){2-3}\cmidrule(lr){4-5}\cmidrule(lr){6-7}
$K$ & Observed & Indep Null & Observed & Indep Null & Observed & Indep Null \\
\midrule
2 & 0.760 & 0.500 & 0.607 & 0.333 & 0.531 & 0.000 \\
3 & 0.548 & 0.333 & 0.476 & 0.200 & 0.482 & 0.000 \\
4 & 0.424 & 0.250 & 0.410 & 0.143 & 0.461 & 0.000 \\
5 & 0.350 & 0.200 & 0.374 & 0.111 & 0.456 & 0.000 \\
\bottomrule
\end{tabular}
\par\medskip
\begin{minipage}{\linewidth}
\footnotesize
\textit{Notes:} Observed values are means across four assets, GMM and HMM models, 20 seeds, and all admissible $(A, B)$ representation pairs at window $=252$, step $=21$. For each pair and window, per-state CVaR ($\alpha=0.05$) is computed from in-state returns, each state sequence is relabelled so that label 0 corresponds to the worst-CVaR state, then the two relabelled sequences are compared pointwise. Independent null expectations are given analytically ($1/K$, $1/(2K-1)$, $0$) and confirmed by Monte Carlo to within $\pm 0.01$ at $n_\text{perm}=500$ per asset. The observed-to-null excess is between $1.6\times$ and $3.4\times$ across $K$, genuinely above chance, but the absolute magnitude is comparable to the partition-level cross-representation ARI in Table~\ref{tab:baseline} ($0.34$--$0.43$). Computation script: \texttt{src/experiments/posthoc\_rank\_aligned\_ordering.py}; outputs to \texttt{outputs/rank\_aligned\_ordering\_summary.csv}.

\textit{Source:} Authors' calculations.
\end{minipage}
\end{table}

\section{Per-pair Permutation Test Details}
\label{app:perm_details}

Table~\ref{tab:perm_per_asset} reports the per-pair distribution at the baseline specification under two complementary nulls and after multiple-comparison correction.

The iid permutation null (random shuffle of one representation's label sequence) is the standard reference but understates dependence: rolling-window labels are persistent, so iid resampling produces null ARI distributions that are too tight, biasing raw $p$-values downward. We therefore additionally report a circular block bootstrap null \citep{PolitisRomano1992CBB} on the same per-pair data, with the block length selected per pair via the Patton-Politis-White (2009) optimum \citep{PattonPolitisWhite2009}. The selected block length captures the within-series autocorrelation of regime labels (per-pair median ${\sim}26$--$27$ trading days, range $[2, 48]$ across pairs and assets at the baseline specification), so the CBB null reproduces the persistence structure of the observed sequence under random alignment. The CBB null is materially wider than iid: aggregate $p$-values rise by approximately a factor of three to four (perm: $0.016$--$0.021$; CBB: $0.063$--$0.071$), and the BH-adjusted survival fraction at $q=0.05$ drops from $97.9\%$--$98.5\%$ under iid to $67.7\%$--$71.2\%$ under CBB. The CBB rate is the more honest of the two; we report both for transparency. We control the false discovery rate via the Benjamini-Hochberg (1995) step-up procedure at $q=0.05$ \citep{BenjaminiHochberg1995}; per-asset BH-survival fractions are reported separately for each null in Table~\ref{tab:perm_per_asset}. Bonferroni control at $n=2{,}000$ pairs requires raw $p<2.5\times 10^{-5}$, below the smallest attainable $p$-value at $B=99$ replications $(1/(B+1)=0.01)$, and is therefore uninformative in this design; BH-FDR is the appropriate correction given that the principal claim is the aggregate magnitude of cross-representation ARI rather than the discovery of any specific pair.

% source: outputs/{BTC-USD,GLD,GSPC,IEF}/step_21/key_results.csv  (per-asset perm + CBB + BH cols)
\begin{table}[!htbp]
\centering
\small
\setlength{\tabcolsep}{4pt}
\caption{Per-pair test of cross-representation ARI by asset, baseline specification (window $=252$, step $=21$, $K=3$, 20 seeds, 7--8 representations). Two nulls: iid permutation and \citet{PolitisRomano1992CBB} circular block bootstrap (CBB) with Patton-Politis-White (2009) block length; $B=99$ replications and 2{,}000 pairs per asset under each. Discovery columns report the Benjamini-Hochberg (1995) step-up fraction at $q=0.05$.}
\label{tab:perm_per_asset}
\begin{tabular}{lcccccc}
\toprule
Asset & Mean ARI & Median $p_\text{perm}$ & \% BH $q<0.05$ (perm) & Median $p_\text{CBB}$ & \% BH $q<0.05$ (CBB) & Median block (days) \\
\midrule
BTC-USD & 0.396 & 0.010 & 98.5 & 0.010 & 71.2 & 26 \\
GLD & 0.389 & 0.010 & 98.2 & 0.010 & 69.4 & 27 \\
GSPC & 0.367 & 0.010 & 98.3 & 0.010 & 70.3 & 26 \\
IEF & 0.367 & 0.010 & 97.9 & 0.020 & 67.7 & 26 \\
\bottomrule
\end{tabular}
\medskip
\footnotesize
\textit{Notes:} Numeric rows are auto-synchronised from \texttt{outputs/<asset>/step\_21/key\_results.csv} by \texttt{src/workflows/paper\_autofill.py} after each posthoc run; manually edited values are overwritten on the next sync. Pairs not surviving BH at $q=0.05$ correspond to (model, seed, rolling window) combinations in which the observed ARI lies near zero; both the observed statistic and the resampled null concentrate near zero for these pairs, so the test correctly fails to reject. The CBB column's BH-survival fraction is materially lower than the iid-permutation column's, as expected: preserving the within-series autocorrelation widens the null distribution and makes raw $p$-values larger. The principal finding (mean cross-representation ARI well below 0.65) is established by the aggregate ARI distribution rather than by universal per-pair rejection, and survives both nulls.
\end{table}

\section{Stability Metric Details}
\label{app:metrics}

\paragraph{Adjusted Mutual Information (AMI).}
As a robustness check on ARI we compute AMI \citep{VinhEppsBailey2010}, which adjusts for chance but weights the full conditional entropy rather than large-cluster agreement \citep{WarrensvanderHoef2022}. AMI follows the same rank ordering as ARI across all model--asset combinations: cross-representation AMI is below temporal AMI for HMM and above or equal for GMM, and the values remain far below any practically useful threshold. This rules out the possibility that the ARI finding is an artefact of ARI's known sensitivity to large-cluster structure.

\paragraph{Per-pair permutation tests.}
For each pair of representations we draw 2{,}000 representative (model, seed, roll) observations and compute a one-sided p-value under two nulls: an iid permutation of one representation's label sequence, and a circular block bootstrap with block length set per pair by the Patton-Politis-White (2009) selector \citep{PolitisRomano1992CBB,PattonPolitisWhite2009}. The CBB preserves the within-series autocorrelation of regime labels (per-asset median selected block length ${\sim}26$--$27$ trading days) under random alignment with the partner labelling, addressing the iid null's anti-conservatism on persistent rolling-window sequences. Both nulls use $B=99$ and the $(e+1)/(B+1)$ correction. We control the false discovery rate via the Benjamini-Hochberg (1995) step-up procedure at $q=0.05$ \citep{BenjaminiHochberg1995}; the per-asset BH-adjusted survival fraction is reported separately for each null in Table~\ref{tab:perm_per_asset}. The CBB rejection rate is the more honest of the two and is materially lower than the iid rate, but the substantive ranking of cross-representation ARI well below temporal ARI is preserved under both nulls. Pairs that do not survive correspond to (model, seed, rolling window) slices where the observed ARI lies near zero; both the observed and resampled distributions concentrate near zero for these pairs, so the test correctly fails to reject. We view the per-pair test as confirming non-random sign rather than magnitude: the substantive finding is that mean ARI is well below 0.65, not merely that it exceeds zero.

\section{Per-asset Infimum ARI (Worst-case Pair)}
\label{app:infimum}

The empirical diagnostic in equation~(2) is formally an infimum over representation pairs; the paper reports the per-asset \emph{mean} as a summary, which understates the worst case. Table~\ref{tab:infimum} reports the per-asset minimum cross-representation ARI across all admissible $(\Phi_1, \Phi_2)$ pairs within each (asset, model) cell, computed at the baseline specification (window $=252$, step $=21$, $K=3$, averaged across 20 seeds and all rolling windows). The worst-case ARI is materially below the corresponding mean: $0.147$--$0.274$ across (asset, model) cells versus the $0.328$--$0.430$ means in the per-cell decomposition. The infimum lies far below any practical stability threshold; the dominant low-agreement pair is rep\_a\_unscaled vs rep\_b for BTC-USD/GLD/IEF and rep\_a\_unscaled vs rep\_e for \SPX{} (the latter being the only cell where rep\_e participates).

% source: outputs/{BTC-USD,GLD,GSPC,IEF}/step_21/results/rep_stability.json  (per-pair ARI; infimum taken over admissible pairs)
\begin{table}[!htbp]
\centering
\small
\setlength{\tabcolsep}{6pt}
\caption{Per-asset infimum cross-representation ARI at $K=3$, window $=252$, step $=21$, with the worst pair identified.}
\label{tab:infimum}
\begin{tabular}{llcl}
\toprule
Asset & Model & Min ARI & Worst pair \\
\midrule
BTC-USD & GMM & 0.274 & rep\_a\_unscaled vs rep\_b \\
BTC-USD & HMM & 0.236 & rep\_a\_unscaled vs rep\_b \\
GLD     & GMM & 0.238 & rep\_a\_unscaled vs rep\_b \\
GLD     & HMM & 0.214 & rep\_a\_unscaled vs rep\_b \\
\SPX{}  & GMM & 0.147 & rep\_a\_unscaled vs rep\_e \\
\SPX{}  & HMM & 0.158 & rep\_a\_unscaled vs rep\_e \\
\IEF{}  & GMM & 0.219 & rep\_a\_unscaled vs rep\_b \\
\IEF{}  & HMM & 0.201 & rep\_a\_unscaled vs rep\_b \\
\bottomrule
\end{tabular}
\par\medskip
\begin{minipage}{\linewidth}
\footnotesize
\textit{Notes:} For each (asset, model) cell, ARI is averaged across 20 seeds and all rolling windows for every admissible representation pair, and the smallest such pair-mean is reported. The minimum is consistently driven by the unstandardised baseline (rep\_a\_unscaled) paired against either an alternative feature philosophy (rep\_b: realised skewness and stability-of-volatility) or a different information source (rep\_e: VIX implied volatility, \SPX{} only). Computed from \texttt{outputs/<asset>/step\_21/results/rep\_stability.json}. \\
\textit{Source:} Authors' calculations.
\end{minipage}
\end{table}

\section{HMM Convergence Audit}
\label{app:convergence}

For bit-exact reproducibility of HMM seed initialisation, the pipeline sets \texttt{OMP\_NUM\_THREADS}, \texttt{MKL\_NUM\_THREADS}, and \texttt{OPENBLAS\_NUM\_THREADS} to $1$ before any numerical work (the snapshot is recorded in \texttt{outputs/run\_manifest.json}), since \texttt{hmmlearn}'s k-means initialisation can otherwise depend on the BLAS thread count.

\texttt{hmmlearn}'s EM optimiser emits a ``Model is not converging'' warning whenever the final iteration's log-likelihood update is non-positive (or below \texttt{tol}); the published pipeline's \texttt{outputs/run.log} contains $54{,}370$ such warnings, which without context is alarming. We therefore parsed the log and binned the warnings by $|\Delta\ell|$ magnitude (script: \texttt{src/experiments/posthoc\_convergence\_audit.py}; output: \texttt{outputs/hmm\_convergence\_audit.csv}). Across the $712{,}420$ HMM fits in the full pipeline (4 assets $\times$ up to $8$ representations $\times$ $2$ models $\times$ $4$ steps $\times$ $20$ seeds $\times$ rolling windows; including the $K\in\{2,4,5\}$ robustness sweep), the $54{,}370$ warnings amount to $7.6\%$ of fits, with median $|\Delta\ell|=8.0\times 10^{-3}$ and the bulk of the mass at machine-precision levels. At a strict $|\Delta\ell|>10^{-3}$ threshold (a delta below which standard practice considers the EM to have stabilised), $36{,}308$ fits qualify as truly non-converged ($5.1\%$ of pipeline fits). Stricter still, $|\Delta\ell|>10^{-1}$ retains $21{,}409$ fits ($3.0\%$). The empirical findings reported in the main text are stable to dropping non-converged fits: cross-representation and temporal ARI computed on the $\geq 95\%$ of converged fits differ from the headline numbers by less than $0.005$ (untabulated). Truly non-converged fits cluster in the rep\_a\_unscaled column, which is consistent with the known HMM-on-unstandardised-scale fragility discussed in the dimension decomposition; for the seven standardised representations the truly non-converged rate is below $2\%$.

\section{State-Number Sensitivity: Model-Class Discussion}
\label{app:k_sens}

The GMM reversal at $K=4$ (cross-representation ARI $0.415$ exceeds temporal ARI $0.302$; Table~\ref{tab:robustness}) admits a natural explanation. GMM, lacking transition dynamics, assigns observations to components purely on the basis of the input feature distribution at estimation time. As $K$ grows, the mixture becomes finer and more sensitive to which observations fall inside the estimation window: a shift of $s=21$ days can reassign a non-trivial fraction of boundary observations, degrading temporal consistency. Cross-representation agreement, by contrast, is governed by which features and windows define the input space, not by window shifts; as long as representations induce similar coarse partitions of the feature space (which they must for nearby specifications), cross-representation ARI is relatively stable across $K$. The HMM does not display the same flip because its transition dynamics regularise the partition across time, but its $K=4$ values still tell a similar story numerically: HMM cross-representation ARI drops from $0.351$ at $K=3$ to $0.314$ at $K=4$ while HMM temporal ARI drops from $0.421$ to $0.308$, so the two stability dimensions converge rather than diverge: the absence of a flip is a sign-level distinction, not a magnitude one. Both models therefore exhibit the same finer-grained-partitions-are-less-stable mechanism; what differs is whether the cross/temporal ranking happens to cross. The apparent ranking of the two instability sources is consequently a function of the granularity of the state space and the model class, not a structural property of the data.

\section{Synthetic Ground-Truth Experiment: DGP Setup}
\label{app:synth_setup}

The baseline DGP generates $T=2{,}200$ daily observations from a three-state first-order Markov chain with transition persistence $p=0.97$ (probability of staying in the same state). States differ in return drift and volatility: the low-risk state has positive drift ($\mu_1=0.04\%$/day, $\sigma_1=0.5\%$), the medium-risk state is near-flat ($\mu_2=0\%$, $\sigma_2=1.2\%$), and the high-risk state has negative drift ($\mu_3=-0.08\%$, $\sigma_3=2.5\%$). A deterministic drift component ($\alpha=0.6$) modulates the mean slightly to produce realistic-looking price paths. We generate 10 independent DGP seeds and for each seed run 3 independent fit seeds, yielding $30$ realisations per representation. The representation pipeline is applied to four of the eight admissible representations (rep\_a, rep\_a\_unscaled, rep\_b, rep\_c1); the GARCH and VIX specifications are excluded because GARCH fitting on short synthetic samples is prohibitively slow and VIX has no natural synthetic counterpart. Using four representations produces 6 cross-representation pairs per window, sufficient for the qualitative comparison.

\paragraph{Pipeline-level interpretation of the recovery shortfall.}
The $0.26$--$0.28$ recovery ARI is a pipeline-level negative result: the rolling HMM/GMM with risk-based features cannot recover a known Markov partition under this DGP, and we do not interpret real-data partitions as approximations to a true regime structure. A natural decomposition of the shortfall: at $p=0.97$ the expected regime duration is $1/(1{-}p)=33$ business days, while the volatility feature uses a $20$-day rolling window (drawdown / tail features use $60$-day windows), so each transition leaves the feature lagging the regime for roughly $20$--$60$ days out of an average $33$-day regime; over a $252$-day estimation window this lag dominates in-sample observability of the regime boundary, producing systematic mis-classification at transitions even when the underlying state structure is information-sufficient in expectation. The synthetic cross-representation ARI ($0.31$--$0.36$ at the baseline $K=3$, $p=0.97$) is moreover of the same order of magnitude as the cross-representation ARI observed on real assets ($0.37$--$0.40$ in Table~\ref{tab:baseline}), which directly rebuts the counter-argument that cross-representation disagreement might be a property of financial-data noise rather than of the rolling HMM/GMM pipeline itself: the disagreement reproduces under a known Markov DGP with i.i.d.\ Gaussian emissions.

\paragraph{Drift ablation.}
To confirm that the recovery failure is not an artefact of the drift component, we repeat the experiment with $\alpha = 0$ (drift removed, states differ only in volatility). On the four-rep specification cited in Section~\ref{sec:discussion}, recovery ARI moves from HMM $0.26$ / GMM $0.28$ at $\alpha=0.6$ to HMM $0.27$ / GMM $0.28$ at $\alpha=0$ (source: \texttt{outputs/synthetic\_groundtruth\_alpha0.csv}), a shift of at most $0.01$. The same invariance holds under the two-rep sanity-check setup of Table~\ref{tab:synth_robustness}. The inability to recover the true partition is therefore driven by the volatility-dominated feature space, not by the small state-specific drift term.

\paragraph{Persistence sweep.}
To confirm that the recovery shortfall is not an artefact of the single chosen persistence level, we sweep $p \in \{0.95, 0.97, 0.99\}$ holding $K=3$ and $T=2{,}200$ fixed (Table~\ref{tab:synth_psweep}). Recovery ARI peaks at $p=0.97$ and falls at both endpoints; cross-representation ARI is essentially flat. The recovery shortfall persists across all three persistence levels, and at $p=0.95$ the gap widens: cross-rep ARI ($0.275$/$0.307$) exceeds recovery ARI ($0.147$/$0.168$), indicating that representations agree with each other more than with truth, consistent with shared misspecification rather than convergence toward the true partition.

\paragraph{State-count sweep.}
To confirm robustness to the chosen state count, we sweep $K \in \{2, 3, 4\}$ holding $p=0.97$ and $T=2{,}200$ fixed (Table~\ref{tab:synth_ksweep}). The recovery shortfall persists across all three $K$ values: recovery ARI ranges $0.16$--$0.27$ (HMM) and $0.18$--$0.29$ (GMM), and cross-representation ARI ranges $0.29$--$0.31$ (HMM) and $0.36$--$0.41$ (GMM), all well below the $0.65$ threshold. Recovery is hardest at $K=2$ (the coarsest partition is also the least informative under volatility-only features); cross-rep ARI is highest at $K=2$ for GMM ($0.407$), reflecting the same coarse-partition agreement, but truth-recovery still fails. At $K=4$ recovery is comparable to $K=3$ for GMM and slightly lower for HMM, consistent with the finer-partition fragility documented in Section~\ref{app:k_sens}.

% source: outputs/synthetic_groundtruth.csv                (full per-seed log)
%         outputs/synthetic_sanity/recovery_table.csv      (alpha=0 vs 0.6 means; per Notes)
\begin{table}[!htbp]
\centering
\small
\setlength{\tabcolsep}{4pt}
\caption{Synthetic experiment recovery ARI under the baseline DGP, with and without the deterministic drift component $\alpha$.}
\label{tab:synth_robustness}
\begin{tabular}{lcc}
\toprule
Drift $\alpha$ & GMM recovery ARI & HMM recovery ARI \\
\midrule
0.0 & 0.209 $\pm$ 0.005 & 0.201 $\pm$ 0.030 \\
0.6 & 0.208 $\pm$ 0.002 & 0.204 $\pm$ 0.028 \\
\bottomrule
\end{tabular}
\par\medskip
\footnotesize
\textit{Notes:} Recovery ARI = Adjusted Rand Index between inferred states and the true DGP partition, across 5 DGP seeds (mean $\pm$ 95\% CI). $p=0.97, K=3, T=2{,}200$ in both rows; two synthetic representations (rep\_u: shorter-horizon rolling features, 20d vol / 60d drawdown / 60d tail; rep\_l: longer-horizon rolling features, 60d vol / 120d drawdown / 120d tail); from \texttt{outputs/synthetic\_sanity/recovery\_table.csv}. \textit{Source:} Authors' calculations.
\end{table}

% source: outputs/synthetic_groundtruth_psweep_summary.csv  (per-p means)
%         outputs/synthetic_groundtruth_psweep.csv          (per-seed inputs)
\begin{table}[!htbp]
\centering
\small
\setlength{\tabcolsep}{4pt}
\caption{Persistence sweep: recovery and cross-representation ARI under the synthetic DGP for $p \in \{0.95, 0.97, 0.99\}$ ($K=3$, $T=2{,}200$, $\alpha=0.6$, 4 representations).}
\label{tab:synth_psweep}
\begin{tabular}{lcccc}
\toprule
& \multicolumn{2}{c}{HMM} & \multicolumn{2}{c}{GMM} \\
\cmidrule(lr){2-3} \cmidrule(lr){4-5}
Persistence $p$ & Recovery ARI & Cross-rep ARI & Recovery ARI & Cross-rep ARI \\
\midrule
0.95 & 0.147 $\pm$ 0.013 & 0.275 $\pm$ 0.014 & 0.168 $\pm$ 0.013 & 0.307 $\pm$ 0.013 \\
0.97 & 0.264 $\pm$ 0.019 & 0.305 $\pm$ 0.020 & 0.276 $\pm$ 0.018 & 0.355 $\pm$ 0.018 \\
0.99 & 0.211 $\pm$ 0.019 & 0.319 $\pm$ 0.018 & 0.265 $\pm$ 0.019 & 0.346 $\pm$ 0.016 \\
\bottomrule
\end{tabular}
\par\medskip
\footnotesize
\textit{Notes:} Each cell is mean $\pm$ standard error across 10 DGP seeds $\times$ 3 fit seeds $\times$ 4 representations ($n=120$). Recovery ARI = ARI between inferred state assignments and the ground-truth Markov chain. Cross-rep ARI = mean ARI across the 6 cross-representation pairs computed within the same fit seed. Both ARIs remain well below the $0.65$ partition-recovery threshold at every persistence level; at $p=0.95$ recovery ARI falls below cross-rep ARI, which is consistent with shared misspecification rather than convergence to truth. From \texttt{outputs/synthetic\_groundtruth\_psweep\_summary.csv}, generated by \texttt{python run.py posthoc\_synthetic\_psweep}. \textit{Source:} Authors' calculations.
\end{table}

% source: outputs/synthetic_groundtruth_ksweep_summary.csv  (per-K means)
%         outputs/synthetic_groundtruth_ksweep.csv          (per-seed inputs)
\begin{table}[!htbp]
\centering
\small
\setlength{\tabcolsep}{4pt}
\caption{State-count sweep: recovery and cross-representation ARI under the synthetic DGP for $K \in \{2, 3, 4\}$ ($p=0.97$, $T=2{,}200$, $\alpha=0.6$, 4 representations).}
\label{tab:synth_ksweep}
\begin{tabular}{lcccc}
\toprule
& \multicolumn{2}{c}{HMM} & \multicolumn{2}{c}{GMM} \\
\cmidrule(lr){2-3} \cmidrule(lr){4-5}
State count $K$ & Recovery ARI & Cross-rep ARI & Recovery ARI & Cross-rep ARI \\
\midrule
2 & 0.158 $\pm$ 0.016 & 0.294 $\pm$ 0.021 & 0.184 $\pm$ 0.015 & 0.407 $\pm$ 0.023 \\
3 & 0.265 $\pm$ 0.019 & 0.305 $\pm$ 0.020 & 0.276 $\pm$ 0.018 & 0.355 $\pm$ 0.018 \\
4 & 0.234 $\pm$ 0.016 & 0.307 $\pm$ 0.014 & 0.287 $\pm$ 0.015 & 0.351 $\pm$ 0.014 \\
\bottomrule
\end{tabular}
\par\medskip
\footnotesize
\textit{Notes:} Each cell is mean $\pm$ standard error across 10 DGP seeds $\times$ 3 fit seeds $\times$ 4 representations ($n=120$). Volatility levels are log-spaced from $0.006$ to $0.024$ across the $K$ states (matching the baseline $K=3$ tuple $(0.006, 0.012, 0.024)$ exactly), so the DGP is comparable across $K$ values up to the partition granularity. The recovery shortfall persists at every $K$: both ARIs remain well below the $0.65$ partition-recovery threshold. From \texttt{outputs/synthetic\_groundtruth\_ksweep\_summary.csv}, generated by \texttt{python run.py posthoc\_synthetic\_ksweep}. \textit{Source:} Authors' calculations.
\end{table}

\section{Supplementary Figures}
\label{app:figures}

\begin{figure}[H]
\centering
\includegraphics[width=1\linewidth]{fig_representation_conflict.png}
\caption{Representation conflict under COVID-19 stress for the \SPX{} (February--June 2020). rep\_a uses volatility (20d), drawdown (60d), MaxDD, VaR/CVaR (60d) with rolling $z$-score standardization; rep\_c1 uses volatility (30d), drawdown (90d), VaR/CVaR (90d) with the same standardization. Both are risk-based representations differing only in window parameterization. The upper bands show state assignments under each representation (HMM, $K=3$; GMM produces qualitatively similar patterns); the bottom strip marks dates where the two representations assign different risk labels. Disagreement clusters in the crisis interval when regime classification matters most. Figure~\ref{fig:disagree_ts} (main text) quantifies disagreement across the full sample.}
\label{fig:rep_conflict}
\end{figure}

\section{Representation-Dimension Decomposition}
\label{app:decomp}

Table~\ref{tab:decomp} reports mean cross-representation ARI broken down by the dimension along which each pair varies (Section~3.2 and~4.1). Pairs are classified into six mutually exclusive categories. ``Standardization'' = rep\_a vs rep\_a\_unscaled (same features and windows, only preprocessing differs). ``Feature subset'' = pairs within \{rep\_a, rep\_c2, rep\_c3\} (same windows, MaxDD or VaR toggle). ``Window length'' = rep\_c1 vs rep\_c3 (identical features, different window parameters). ``Mixed window + feature'' = rep\_a vs rep\_c1 etc.\ (simultaneous window and feature change). ``Volatility estimator'' = rep\_d vs any rolling-window rep (GARCH vs rolling realized). ``Feature philosophy'' = rep\_b vs vol/dd family (realized skewness and vol-of-vol vs risk measures). ``VIX signal'' = rep\_e vs any other rep (\SPX{} only; backward-looking realised vs forward-looking implied).

% source: outputs/repr_decomp_summary.csv     (mean ARI per variation dimension)
%         outputs/repr_variance_decomp.csv    (eta^2 decomposition referenced in main text)
\begin{table}[!htbp]
\centering
\small
\setlength{\tabcolsep}{5pt}
\caption{Mean cross-representation ARI by variation dimension (mean $\pm$ 95\% CI across all relevant pairs, averaged over assets and models unless noted).}
\label{tab:decomp}
\begin{tabular}{lccc}
\toprule
Variation dimension & Pairs & Mean ARI & 95\% CI \\
\midrule
Feature subset (same windows)           & 3 & 0.551 & $\pm 0.039$ \\
Volatility estimator (GARCH vs rolling) & 5 & 0.454 & $\pm 0.026$ \\
Window length only                      & 1 & 0.352 & $\pm 0.022$ \\
Feature philosophy (skew/stab vs risk)  & 6 & 0.328 & $\pm 0.022$ \\
Standardization only                    & 1 & 0.307 & $\pm 0.022$ \\
Mixed (window + feature)                & 3 & 0.306 & $\pm 0.017$ \\
VIX signal (\SPX{} only)                & 7 & 0.308 & $\pm 0.032$ \\
\midrule
All pairs (baseline)                    & 21+7 & \multicolumn{2}{c}{see Table~\ref{tab:baseline}} \\
\bottomrule
\end{tabular}
\par\medskip
\begin{minipage}{\linewidth}
\footnotesize
\textit{Notes:} Pair counts for the 7-rep universal class; rep\_e adds 7 additional pairs for \SPX{} only. Values sorted by mean ARI descending. Computed from \texttt{repr\_decomp\_summary.csv} via \texttt{src/experiments/posthoc\_repr\_decomp.py}. For the seven universal-class rows, CI = $1.96 \times$SE across the eight per-(asset, model) group means within each category. The VIX row is \SPX{}-only, so we instead report $1.96 \times$SE across the 14 (model, pair) observations within the asset (7 rep\_e-pair $\times$ 2 models); a normal-approximation CI on the two (asset, model) group means alone would be misleadingly tight ($\pm 0.005$) given $n=2$.

\textit{Source:} Authors' calculations.
\end{minipage}
\end{table}

\paragraph{Variance decomposition.}
\label{app:variance_decomp}
To separate the contribution of representation choice from estimation instability, we decompose the cross-seed variance of ARI into between-representation-pair and within-pair (across-seeds) components via $\eta^2 = \text{SS}_{\text{between-pair}} / \text{SS}_{\text{total}}$. Across assets and models, GMM $\eta^2$ averages $0.32$ (range $[0.31, 0.33]$) and HMM $\eta^2$ averages $0.21$ (range $[0.19, 0.22]$). The higher GMM eta-squared is consistent with its narrower seed-level CI in Table~\ref{tab:baseline} ($\pm 0.002$ vs $\pm 0.033$): under GMM, a third of the ARI variance is attributable to which pair is compared rather than which seed is used. Under HMM the share falls to a fifth (the rest is seed-driven), so HMM cross-representation ARI carries more uncertainty around the same representation-driven signal. Underlying numbers in \texttt{outputs/repr\_variance\_decomp.csv}.

\section{Admissible Class as Registered Scope}
\label{app:scope}

\paragraph{Per-dimension justification of the admissible class.}
The eight representations vary along five documentation-relevant dimensions, each tied to a specific entry point in routine model paperwork rather than ad-hoc design: (i)~standardization (rolling $z$-score vs none), the default preprocessing decision required by SR~11-7's data-and-assumptions principle \citep{SR117}; (ii)~feature subset (MaxDD or VaR inclusion), the tail-risk choice most often varied during regime-model recalibration \citep{Guidolin2011,RadLowMiffreFaff2023}; (iii)~rolling horizon (20/60d vs 30/90d), the canonical realised-volatility window choice \citep{AngTimmermann2012,Kirby2023FRL}; (iv)~volatility estimator (rolling realised vs GARCH(1,1)), the central methodological fork in FRTB market-risk modelling \citep{BCBSD457,BlazsekKongShadoff2025}; (v)~information source (backward-looking realised risk vs forward-looking VIX), distinguishing physical from risk-neutral measure. The first six representations share a rolling-window architecture and substantial mutual overlap (five of six use VaR/CVaR as tail features; four of six share $\text{vol}=20$, $\text{dd}=60$, $\text{tail}=60$); this redundancy is deliberate, so ARI within near-redundant specifications constitutes a conservative lower bound on the disagreement under a broader class.

\paragraph{What the registered-scope reading implies.}
Conclusions are statements about empirical agreement \emph{within} the admissible class above; they are not universal claims about regime models in the abstract. We frame the contribution as methodological: placing cross-representation agreement on the model-validation agenda, and showing that under one fully-specified, regulator-relevant class the agreement is materially below the partition-recovery threshold. A reader working under a strictly narrower class can read off the relevant cell of Table~\ref{tab:decomp}: the highest within-class mean ARI we observe (feature-subset variation under common windows) is $0.55\pm 0.04$, still below the $0.65$ threshold. A reader working under a broader class (multivariate features, hourly resolution, alternative model families) would by construction obtain lower agreement, because additional admissible variation can only weaken pairwise overlap. The class chosen here is intentionally narrow rather than ambitious; the reported numbers are best read as a conservative lower bound on the disagreement realised in production specification spaces.

\section{Robustness and Alternative Interpretations}
\label{app:dialectical}

The central findings are robust in headline terms but warrant dialectical scrutiny on five points where an alternative reading is plausible.

\emph{(i) Threshold choice.} The 0.65 benchmark from \citet{Steinley2004} was calibrated against simulation studies outside finance. A reader skeptical of that specific cut-off can read Table~\ref{tab:baseline}: cross-representation ARI remains at 0.34--0.43 across all models and assets, so even under a substantially looser 0.50 threshold the partition-level finding stands. AMI, which corrects differently for large clusters, produces the same qualitative ranking.

\emph{(ii) Disentangling non-stationarity from estimation noise at large step.} At step~$=252$ (non-overlapping windows), the disjoint temporal ARI of 0.14 could reflect genuine regime non-persistence, estimation instability across independently fitted models separated by one year, or both; the data cannot separate these. We report the value as a lower bound on temporal stability and do not infer non-persistence from it. The cross-representation ARI at step~$=252$ (0.38) is the more informative quantity for this paper's thesis, since it is computed on identical observations.

\emph{(iii) Ordering nulls and the need for a matching-free metric.} The pipeline-native ordering metric (Hungarian + CVaR) has both nulls (the single-side permutation and the independent-random-partition) empirically in the $0.41$--$0.99$ range across (asset, $K$) rather than the $1/K$ implied by a symmetric chance baseline (Section~\ref{app:ordering}). An earlier draft of this paper claimed ``observed ordering exceeds the independent null across $K\in\{3,4,5\}$''; that claim is retracted, because the independent null is constructively biased upward by the same Hungarian + CVaR matching rule used for the observed value. Under a matching-free rank-aligned metric with a rigorous $1/K$ null (Table~\ref{tab:ordering_rank_aligned}), observed top-1 Jaccard at $K=3$ is $0.47$ (vs null $0.20$), an honest $2.4\times$ excess but an order of magnitude smaller than the $0.93$ top-1 reported by the biased metric. The governance recommendation in Section~\ref{sec:conclusion} is therefore based on this matching-free metric, not the Hungarian one, and treats ordering agreement at the same level of scepticism as partition-level agreement.

\emph{(iv) The $\eta^2$ contrast between GMM and HMM.} GMM $\eta^2 \approx 0.32$ versus HMM $\eta^2 \approx 0.21$ does \emph{not} imply that representations matter more for GMM in an absolute sense. The ratio is higher under GMM because GMM is more numerically deterministic (lower seed-level variance in the denominator); the between-pair variance is similar for both models ($0.016$--$0.018$ for GMM vs $0.010$--$0.012$ for HMM). The correct reading is that HMM ARI carries more estimation noise around the same representation-driven signal.

\emph{(v) Synthetic recovery is \emph{suggestive}, not conclusive.} The inability of the pipeline to recover the true partition under a stylized DGP (recovery ARI $\approx 0.15$--$0.27$ across persistence levels) is consistent with instability, but we caution against over-interpreting it for three reasons. First, the DGP has expected regime duration $1/(1{-}p)$ business days while the volatility feature uses a 20-day rolling window; measurement lag alone introduces mis-classification at transitions, which is a finite-sample observability issue rather than a representation-invariance failure. Second, the synthetic class is limited to four of the eight representations (GARCH and VIX are excluded), so class coverage is narrower than the main experiments. Third, the recovery shortfall is robust across both DGP-parameter sweeps reported in Section~\ref{app:synth_setup}: persistence $p \in \{0.95, 0.97, 0.99\}$ at fixed $K=3$ (Table~\ref{tab:synth_psweep}) and state count $K \in \{2, 3, 4\}$ at fixed $p=0.97$ (Table~\ref{tab:synth_ksweep}). At $p=0.95$ recovery ARI even falls below cross-representation ARI, which strengthens rather than weakens the instability reading. The experiment therefore \emph{weakens} the irreducible-plurality interpretation but does not rule it out, and the governance argument we make in Section~\ref{sec:conclusion} is independent of whether instability or plurality is the causal mechanism.

\emph{What would change our conclusion?} The partition-level finding (cross-rep ARI $\ll 0.65$) would be overturned if (a) the admissible class were restricted to specifications that share the same volatility horizon, feature set, and standardization (Table~\ref{tab:decomp} shows feature-subset pairs at $0.55 \pm 0.04$, still below but within roughly one threshold-gap of the $0.65$ poor-recovery cutoff), or (b) a single representation could be shown to dominate the class on a governance-relevant criterion. Neither is a weakness of the diagnostic itself: restricting the class amounts to declaring one representation canonical, which is precisely the audit failure the diagnostic is designed to surface.

\section{Worked Numeric Example: \SPX{} 2020-Q1}
\label{app:worked_example}

To make the day-to-day workflow described in Section~\ref{sec:conclusion} concrete, this appendix walks through the computation on a single asset and a single date. The example uses \SPX{} on 2020-03-23 (the COVID-19 trough), where the cross-representation CVaR envelope reaches its in-sample peak.

\paragraph{Step 1: Inputs.}
The model owner registers six admissible representations from Table~\ref{tab:rep_definitions} (rep\_a, rep\_a\_unscaled, rep\_b, rep\_c1, rep\_c2, rep\_c3), each fitted to \SPX{} adjusted-close prices through 2020-03-23 with $K=3$ HMM and a $W=252$ trailing window. Each representation produces a hard state assignment $z_t^\Phi \in \{0,1,2\}$ on every trading day in the window.

\paragraph{Step 2: Per-day cross-representation labels and conditional CVaR.}
On 2020-03-23 the six universally-applied representations split into two clusters by their hard state assignment for \SPX{}: rep\_a, rep\_a\_unscaled, rep\_c2, and rep\_c3 assign \SPX{} to their highest-volatility state (label 2 after volatility ordering), while rep\_b and rep\_c1 assign it to the medium-volatility state (label 1). The regime-conditional 5\% CVaR (computed from in-state log returns within the trailing 252-day window) therefore differs in kind across the two clusters: representations that assign the high-volatility label produce a conditional loss of order $-700$~bp/day, while representations that assign the medium-volatility label produce a conditional loss of order $-200$~bp/day. The unconditional 5\% CVaR for \SPX{} on the same trailing window is approximately $-230$~bp/day. Exact per-representation values for this date are produced by \texttt{outputs/var\_spread\_summary.csv} when the pipeline is rerun against the registered data vintage; we report the structural pattern here rather than 4-decimal point estimates because the latter are sensitive to seed and trailing-window endpoint.

\paragraph{Step 3: Envelope and trigger.}
The cross-representation spread is $S_t = \max_\Phi \widehat{\mathrm{CVaR}}_t^\Phi - \min_\Phi \widehat{\mathrm{CVaR}}_t^\Phi$, which on this date is approximately $500$~bp, about $200\%$ of the unconditional CVaR. Aggregated over the 2020-03-15 to 2020-04-15 window, the per-day mean spread reaches $267\%$ (Table~\ref{tab:cvar_spread}, \SPX{} GMM COVID-19 peak). The materiality trigger is $S_t / |\mathrm{CVaR}_\text{unc}| > 0.10$, which is breached with two orders of magnitude of headroom on this day. The cross-representation ARI on the same trailing window is well below the $0.65$ trigger.

\paragraph{Step 4: Action.}
Both triggers fire simultaneously (envelope $\gg 10\%$, ARI $< 0.65$), so the model-validation workflow escalates the day's regime classification to MRM committee review. Two responses are admissible under the recommendation: post the conservative end of the envelope (the more-negative cluster, of order $-700$~bp) for capital purposes, or report a mean-and-range envelope spanning the medium-volatility cluster ($\sim$$-200$~bp) to the high-volatility cluster ($\sim$$-700$~bp) and document the representation uncertainty as a known model limitation under SR~11-7. The point of the worked example is that both responses are documentable, repeatable, and timestamped; the workflow does not require the modeller to choose a ``correct'' representation, only to disclose the disagreement.

\paragraph{Reproducibility.}
The complete reproducibility capsule (data + code + locked dependency versions) is registered as a Zenodo deposit; the DOI is added on acceptance.

\bibliography{refs}

\end{document}
